\section{GALOIS system architecture}
Our implementation orchestrates the interaction between a declarative SQL interface and the probabilistic nature of LLMs through four distinct stages: parsing, logical planning, physical execution, and local post-processing. This section details the core components implemented in the \texttt{Galois} orchestration class.


\subsection{Orchestration and Schema Management}

The entry point of the system is the \texttt{Galois} class, which initializes the workflow by parsing the input SQL query into an Abstract Syntax Tree (AST) using the \texttt{sql\_query\_parser} module. Unlike standard text-to-SQL approaches, GALOIS does not immediately prompt the LLM with the raw query, instead it inspects the database structure using the \texttt{GaloisSchemaManager}.
This component connects to a read-only DuckDB instance to retrieve accurate metadata (column names, primary keys, and data types). This schema awareness allows the system to construct semantically valid JSON schemas for the prompts, ensuring that the LLM's output adheres to the strict typing required for the subsequent relational operations.

\subsection{The Execution Engine}
The retrieval logic is encapsulated within the \texttt{GaloisExecutor} class. This component abstracts the complexity of the LLM interaction (via LangChain) and implements the iterative retrieval loops.

\textbf{State Management and Deduplication}: To handle the context window limitations of LLMs, the executor utilizes a custom \texttt{GaloisMemory} class. This memory component maintains a history of retrieved tuples and implements a strict deduplication mechanism based on primary key hashing \texttt{(\_get\_key\_values)}. This ensures that subsequent iterations (up to a configurable max\_iter) prompt the model to generate only novel data.

\textbf{Iterative Scan Operators}: The executor supports two physical operators:
\begin{itemize}
    \item \textbf{Table-Scan}: Iteratively requests complete tuples (keys and attributes) until convergence or iteration limits are reached.
    \item \textbf{Key-Scan}: Decomposes retrieval into two phases. First, it iterates to collect all unique primary keys. Once the keys are consolidated, it issues a final "Tuple Request"  to retrieve non-key attributes for the collected identifiers.
\end{itemize}

\subsection{}




